%%% The main file. It contains definitions of basic parameters and includes all other parts.

% Meta-data of your thesis (please edit)
\input metadata.tex

% Generate metadata in XMP format for use by the pdfx package
\input xmp.tex

%% Settings for single-side (simplex) printing
% Margins: left 40mm, right 25mm, top and bottom 25mm
% (but beware, LaTeX adds 1in implicitly)
% \documentclass[12pt,a4paper]{report}
% \setlength\textwidth{145mm}
% \setlength\textheight{247mm}
% \setlength\oddsidemargin{15mm}
% \setlength\evensidemargin{15mm}
% \setlength\topmargin{0mm}
% \setlength\headsep{0mm}
% \setlength\headheight{0mm}
% \openright makes the following text appear on a right-hand page
% \let\openright=\clearpage

%% Settings for two-sided (duplex) printing
% \documentclass[12pt,a4paper,twoside,openright]{report}
% \setlength\textwidth{145mm}
% \setlength\textheight{247mm}
% \setlength\oddsidemargin{14.2mm}
% \setlength\evensidemargin{0mm}
% \setlength\topmargin{0mm}
% \setlength\headsep{0mm}
% \setlength\headheight{0mm}
% \let\openright=\cleardoublepage

%% If the thesis has no printed version, symmetric margins look better
\documentclass[12pt,a4paper]{report}
\setlength\textwidth{145mm}
\setlength\textheight{247mm}
\setlength\oddsidemargin{10mm}
\setlength\evensidemargin{10mm}
\setlength\topmargin{0mm}
\setlength\headsep{0mm}
\setlength\headheight{0mm}
\let\openright=\clearpage

%% Generate PDF/A-2u
\usepackage[a-2u]{pdfx}

%% Prefer Latin Modern fonts
\usepackage{lmodern}
% If we are not using LuaTeX, we need to set up character encoding:
\usepackage{iftex}
\ifpdftex
\usepackage[utf8]{inputenc}
\usepackage[T1]{fontenc}
\usepackage{textcomp}
\fi

%% Further useful packages (included in most LaTeX distributions)
\usepackage{amsmath}        % extensions for typesetting of math
\usepackage{amsfonts}       % math fonts
\usepackage{amsthm}         % theorems, definitions, etc.
\usepackage{bm}             % boldface symbols (\bm)
\usepackage{booktabs}       % improved horizontal lines in tables
\usepackage{caption}        % custom captions of floating objects
\usepackage{dcolumn}        % improved alignment of table columns
\usepackage{floatrow}       % custom float environments
\usepackage{graphicx}       % embedding of pictures
\usepackage{indentfirst}    % indent the first paragraph of a chapter
\usepackage[nopatch=item]{microtype}   % micro-typographic refinement
\usepackage{paralist}       % improved enumerate and itemize
\usepackage[nottoc]{tocbibind} % makes sure that bibliography and the lists
			    % of figures/tables are included in the table
			    % of contents
\usepackage{xcolor}         % typesetting in color

% The hyperref package for clickable links in PDF and also for storing
% metadata to PDF (including the table of contents).
% Most settings are pre-set by the pdfx package.
\hypersetup{unicode}
\hypersetup{breaklinks=true}

% Packages for computer science theses
\usepackage{algpseudocode}  % part of algorithmicx package
\usepackage{algorithm}
\usepackage{fancyvrb}       % improved verbatim environment
\usepackage{listings}       % pretty-printer of source code

% You might want to use cleveref for references
% \usepackage{cleveref}

% Set up formatting of bibliography (references to literature)
% Details can be adjusted in macros.tex.
%
% BEWARE: Different fields of research and different university departments
% have their own customs regarding bibliography. Consult the bibliography
% format with your supervisor.
%
% The basic format according to the ISO 690 standard with numbered references
\usepackage[natbib,style=iso-numeric,sorting=none]{biblatex}
% ISO 690 with alphanumeric references (abbreviations of authors' names)
%\usepackage[natbib,style=iso-alphabetic]{biblatex}
% ISO 690 with references Author (year)
%\usepackage[natbib,style=iso-authoryear]{biblatex}
%
% Some fields of research prefer a simple format with numbered references
% (sorting=none tells that bibliography should be listed in citation order)
%\usepackage[natbib,style=numeric,sorting=none]{biblatex}
% Numbered references, but [1,2,3,4,5] is compressed to [1-5]
%\usepackage[natbib,style=numeric-comp,sorting=none]{biblatex}
% A simple format with alphanumeric references:
%\usepackage[natbib,style=alphabetic]{biblatex}

% Load the file with bibliography entries
\addbibresource{bibliography.bib}

% Our definitions of macros (see description inside)
\input macros.tex

%%% Title page and various mandatory informational pages
\begin{document}
%%% Title page of the thesis and other mandatory pages

%%% Inscriptions at the opening page of the hard cover

% We usually do not typeset the hard cover, but if you want to do it, change \iffalse to \iftrue
\iffalse

\pagestyle{empty}
\hypersetup{pageanchor=false}
\begin{center}

\large
Charles University

\medskip

Faculty of Mathematics and Physics

\vfill

{\huge\bf\ThesisTypeTitle}

\vfill

{\huge\bf\ThesisTitle\par}

\vfill
\vfill

\hbox to \hsize{\YearSubmitted\hfil \ThesisAuthor}

\end{center}

\newpage\openright
\setcounter{page}{1}

\fi

%%% Title page of the thesis

\pagestyle{empty}
\hypersetup{pageanchor=false}
\begin{center}

\centerline{\mbox{\includegraphics[width=166mm]{img/logo-en.pdf}}}

\vspace{-8mm}
\vfill

{\bf\Large\ThesisTypeTitle}

\vfill

{\LARGE\ThesisAuthor}

\vspace{15mm}

{\LARGE\bfseries\ThesisTitle\par}

\vfill

\Department

\vfill

{
\centerline{\vbox{\halign{\hbox to 0.45\hsize{\hfil #}&\hskip 0.5em\parbox[t]{0.45\hsize}{\raggedright #}\cr
Supervisor of the \ThesisTypeName{} thesis:&\Supervisor \cr
\ifx\ThesisType\TypeRig\else
\noalign{\vspace{2mm}}
Study programme:&\StudyProgramme \cr
\fi
}}}}

\vfill

Prague \YearSubmitted

\end{center}

\newpage

%%% A page with a solemn declaration to the thesis

\openright
\hypersetup{pageanchor=true}
\vglue 0pt plus 1fill

\noindent
I declare that I carried out this \ThesisTypeName{} thesis on my own, and only with the cited
sources, literature and other professional sources.
I understand that my work relates to the rights and obligations under the Act No.~121/2000 Sb.,
the Copyright Act, as amended, in particular the fact that the Charles
University has the right to conclude a license agreement on the use of this
work as a school work pursuant to Section 60 subsection 1 of the Copyright~Act.

\vspace{10mm}

\hbox{\hbox to 0.5\hsize{%
In \hbox to 6em{\dotfill} date \hbox to 6em{\dotfill}
\hss}\hbox to 0.5\hsize{\dotfill\quad}}
\smallskip
\hbox{\hbox to 0.5\hsize{}\hbox to 0.5\hsize{\hfil Author's signature\hfil}}

\vspace{20mm}
\newpage

%%% Dedication

\openright

\noindent
\Dedication

\newpage

%%% Mandatory information page of the thesis

\openright
{\InfoPageFont

\vtop to 0.5\vsize{
\setlength\parindent{0mm}
\setlength\parskip{5mm}

Title:
\ThesisTitle

Author:
\ThesisAuthor

\DeptType:
\Department

Supervisor:
\Supervisor, \SupervisorsDepartment

Abstract:
\Abstract

Keywords:
{\def\sep{\unskip, }\ThesisKeywords}

\vfil
}

% In Czech study programmes, it is mandatory to include Czech meta-data:

\ifx\StudyLanguage\LangCS

\vtop to 0.49\vsize{
\setlength\parindent{0mm}
\setlength\parskip{5mm}

Název práce:
\ThesisTitleCS

Autor:
\ThesisAuthor

\DeptTypeCS:
\DepartmentCS

Vedoucí bakalářské práce:
\Supervisor, \SupervisorsDepartment

Abstrakt:
\AbstractCS

Klíčová slova:
{\def\sep{\unskip, }\ThesisKeywordsCS}

\vfil
}

\fi

}

\newpage

%%% Further pages will be numbered, which LaTeX automatically enables at the next chapter start


%%% A page with automatically generated table of contents of the thesis

\tableofcontents

%%% Each chapter is kept in a separate file
\chapter{Introduction}

\noindent In computer science and software engineering, the idea of software verification is becoming more and more important. Verification checks whether the program or system operates correctly and fulfills the given properties. The goal is to detect bugs and errors during the development process.

One of the frameworks that is becoming popular is the Constrained Horn Clauses(CHC)\cite{10.1007/978-3-031-13185-1_2} framework. CHC is a fragment of First Order Logic modulo constraints that captures many program verification problems as constraint solving. The main advantage of CHC is that it separates modeling from solving by translating the program's behavior and properties into constrained language and then using a specialized CHC solver to solve various verification tasks across programming languages by deciding the satisfiability problem of a CHC system.

Golem\cite{blicha_golem_2023} is one such solver, which integrates the interpolating SMT solver OpenSMT\cite{10.1007/978-3-319-40970-2_35}. Golem currently implements six model-checking algorithms to solve the CHC satisfiability problem.

On top of solving the CHC satisfiability problem, each engine in Golem provides a validity witness for their answer. In software verification we can think of these witnesses as invariants for SAFE answers and a path to a counterexample for UNSAFE answer. By providing these witnesses, we can ensure that the engines answer is correct. Also to check, wheter the engine doesn't give false witnesses, Golem has built in an internal validator to verify the correctness of the witness. 

One such algorithm, that can be used to solve these problems is PDKind (Property-Directed K-induction)\cite{7886665}. PDKind is an IC3\cite{6148908}-based algorithm, that separates reachability checking and induction reasoning, allowing to replace the induction core with k-induction.

\section{Goals}
\noindent The goal of this work is to create a PDKind engine and integrate it into the Golem Horn Solver\cite{blicha_golem_2023}. 

We will start with defining the background around SMT solving. Then, we will introduce the Golem solver and analyze the PDKind algorithm to design a solution that integrates the PDKind algorithm as an engine into the Golem solver. Afterward, we will implement this design and add the generation of validity witnesses to the implementation. We will test the correctness and efficiency of the implementation using a set of benchmarks and compare the results with other engines. The correctness of the witnesses will be verified using a set of tests that utilize an internal function of the Golem solver, which can verify the correctness of the witnesses.


\chapter{PDKind}

% Describe the individual procedures of the algorithm and in each part analyze how the implementation went. If there were more ways to do it, compare them.
% The last section should describe how we also added validity witnesses because Golem and other engines in it produce them).
\noindent PDKind is an IC3\cite{6148908}-based algorithm, that separates reachability checking and induction reasoning, allowing the induction core to be replaced with k-induction.

In this chapter, we describe and analyze the PDKind algorithm. In sections 4.1-4.3, we will be using methodology, algorithms, and definitions described in \cite{7886665}.
After reading this chapter, the reader should be able to understand the parts of the PDKind algorithm and orientate in our implementation.
%We will also describe the modifications of the algorithm and the decisions we made when implementing some abstract data structures from the paper.




\section{Reachability checking procedure}\label{ReachabilityChecking}

\noindent To understand the following, we need to introduce a few definitions.

\vspace{\baselineskip}\noindent Let us have a formula in the form

\begin{equation}
    A(\vec{x}) \wedge T[B]^k(\vec{x},\vec{y},\vec{w}) \wedge C(\vec{y})
\end{equation}

\vspace{\baselineskip}\begin{definition}
                          For \( k > 1 \), \( T[F]^k(\vec{x}, \vec{x}') \) is defined as
                          \begin{equation*}
                              T(\vec{x}, \vec{w}_1) \land \bigwedge_{i=1}^{k-2} \left( F(\vec{w}_i) \land T(\vec{w}_i, \vec{w}_{i+1}) \right) \land F(\vec{w}_{k-1}) \land T(\vec{w}_{k-1}, \vec{x}')
                          \end{equation*}

                          \noindent where $\vec{w}$ are state variables in the intermediate states.
\end{definition}


\vspace{\baselineskip}
\begin{definition}{\textbf{(Interpolant)}}
    \cite{7886665} If formula (4.1) is unsatisfiable, then $I(\vec{y})$ is an interpolant if
    \label{interpolation}
    \begin{enumerate}
        \item \( A(\vec{x}) \land T[B]^k(\vec{x}, \vec{w}, \vec{y}) \Rightarrow I(\vec{y}) \), and
        \item \( I(\vec{y}) \) and \( C(\vec{y}) \) are inconsistent.
    \end{enumerate}
\end{definition}

\vspace{\baselineskip}
\begin{definition}{\textbf{(Generalization)}}
    \cite{7886665} If formula (4.1) is satisfiable, then $G(\vec{x})$ is a generalization if
    \label{generalization}
    \begin{enumerate}
        \item $G(\vec{x}) \Rightarrow \exists \vec{y}, \vec{w}$  $T[B]^k(\vec{x}, \vec{y}, \vec{w}) \wedge C(\vec{y})$, and
        \item $G(\vec{x})$ and $A(\vec{x})$ are consistent.
    \end{enumerate}
\end{definition}

\renewcommand{\figurename}{Algorithm}
\begin{figure}[h]
    \begin{mdframed}
        \begin{algorithmic}[1]
            \State \textbf{Input:} Target state $F$, maximum steps $k$
            \State \textbf{Data:} Reachability frames $R$, initial states $I$, transition states $T$
            \State \textbf{Output:} True if $F$ is reachable in $k$ steps, False otherwise

            \If{$k = 0$}
                \State \Return{CheckSAT($I \wedge T^0 \wedge F$)}
            \EndIf

            \While{true}
                \If{CheckSAT($R_{k-1} \wedge T \wedge F$)}
                    \State $G \gets$ Generalize($R_{k-1}, T, F$)
                    \If{Reachable($G, k-1$)}
                        \State \Return{true}
                    \Else
                        \State $E \gets$ Explain($G, k-1$)
                        \State $R_{k-1} \gets R_{k-1} \cup E$
                    \EndIf
                \Else
                    \State \Return{false}
                \EndIf
            \EndWhile

        \end{algorithmic}
    \end{mdframed}
    \caption{Reachable method}\label{alg:1}
\end{figure}

\vspace{\baselineskip}
\noindent Method shown in Algorithm \ref{alg:1} tries to reach the initial states backwards by using a depth-first search strategy.

To check if \( F \) is reachable from the initial states in \( k \) steps, we
first check whether \( F \) is reachable in one transition from the previous
frame \( R_{k-1} \). If there is no such transition, then \( F \) is not
reachable in \( k \) steps. Otherwise, we get a state that satisfies \( R_{k-1}
\) and from which \( F \) is reachable in a single step. 

We then call a
generalization procedure, which produces a formula \( G \) (see Definition
\ref{generalization}). Here, generalization helps avoid redundant checks by creating an abstraction of the counterexamples. Instead of storing the given state, it generates a formula covering multiple possible counterexamples, which reduces the number of times the reachability procedure gets called.

Then, using a DFS strategy, we recursively
check whether \( G \) is reachable from the initial states.
If \( G \) is reachable, then \( F \) is also reachable, and the procedure terminates. Otherwise,
we can learn an explanation through interpolation (see Definition
\ref{interpolation}) and eliminate \( G \) by adding the
explanation into the frame \( R_{k-1} \).


\section{Push procedure} \label{Push}

\begin{definition}{\textbf{(Induction Frame)}}
    \cite{7886665}\label{def:IFrame} A set of tuples $F \subset \mathbb{F} \times \mathbb{F}$, where $\mathbb{F}$ is a set of all state formulas in theory $T$, is an induction frame at index $n$ if $(P, \neg P) \in F$ and $\forall (lemma, counterExample) \in F$:
    \begin{itemize}
        \item $lemma$ is valid up to $n$ steps and refutes $counterExample$
        \item $counterExample$ states can be extended to a counterexample to $P$.
    \end{itemize}
\end{definition}

\vspace{\baselineskip}
\noindent Induction frame tracks the strengthening of invariants over number of iterations. Every frame consists of a formula that has been refined to eliminate an invalid state.

The procedure shown in Algorithm \ref{alg:3} is the core of the PDKind algorithm.
It takes the invariant and tries to construct a stronger inductive proof by analyzing its counterexamples and strengthening constraints, while switching between \\
generalization and reachability checking.
We will break it down into smaller pieces to understand how it works.


The first step (line 11) is to check whether \( lemma \) is
\( k \)-inductive (see Definition \ref{Def:kind}). This is done by checking if the formula $(F_{ABS} \wedge T_k
\wedge \neg lemma)$ is satisfiable. We can notice, that \( CheckSAT \) call returns model $m_1$ which
is an assignment of values to variables that satisfies the formula.
If the formula $(F_{ABS} \wedge T_k \wedge \neg lemma)$ isn't satisfiable, we
can push a new obligation to our new induction frame \( G \) and continue.
Otherwise, we get the model $m_1$ and save it for later.

In the second part (line 16), we check if the \( counterExample \) is
reachable. If it is, we get model $m_2$ and generalize it to
$g_2$. We know, that we can reach \( \neg P \) from $g_2$,
so we need to check if $g_2$ is reachable from the initial states. If it is
reachable, the property is invalid, and we mark \( isInvalid \gets true \).


On line 19, we can see that \( Reachable() \) accepts more arguments and
returns more values than shown in Algorithm
\ref{alg:1}. This new
$Reachable(i,j,F)$ method checks if $F$ is reachable in \( k \) steps
where $i\leq k \leq j$.
If this call wasn't successful, we get an interpolant $i_1$ and assign
it to $g_3$ which eliminates $g_2$. We found a new induction
obligation $(g_3, g_2)$, which is a strengthening of \( F \). Now,
we can try again with a potential counterexample eliminated.

Last step is to analyze the induction failure. From the first check we have a
model $m_1$, which is a counterexample to the k-inductiveness of
\( lemma \). We again get $g_1$ as a generalization of $m_1$
and check if $g_1$ is reachable from initial states. If it is
reachable, we replace \( lemma \) with weaker \( \neg counterExample \)
and push this new obligation to \( F \) and \( G \). On the other hand,
if $g_1$ is not reachable, we strengthen \( lemma \) with
$g_3$ and push this new obligation to \( F \).

\begin{figure}[H]
    \begin{mdframed}
        \begin{algorithmic}[1]
            \State \textbf{Input:} Induction frame $F$, $n$, $k$
            \State \textbf{Output:} Old induction frame $F$, new induction frame $G$, $n_p$, $isInvalid$

            \State push elements of $F$ to queue $Q$
            \State $G \gets$ \{\}
            \State $n_p \gets n + k$
            \State $invalid \gets false$
            \While {$\neg invalid$ and $Q$ is not empty}
                \State $(lemma, counterExample)  \gets$ $Q$.pop()
                \State $F_{ABS} \gets \bigwedge a_i$ ,where ($a_i$, $b_i$) \in $F$ \forall $i$ \in \{1,...,$F$.length()\}
                \State $T_k \gets T[F_{ABS}]^k$ by definition
                \State $(s_1, m_1) \gets$ CheckSAT($F_{ABS}, T_k, \neg lemma$)
                \If{$\neg s_1$}
                    \State $G \gets G \cup (lemma, counterExample)$
                    \State Continue
                \EndIf
                \State $(s_2, m_2) \gets$ CheckSAT($F_{ABS} \wedge T_k \wedge counterExample$)
                \If{$s_2$}
                    \State $ g_2 \gets $ Generalize($m_2, T_k, counterExample$)
                    \State $(r_1, i_1, n_1) \gets $ Reachable($n-k+1$, $n$, $g_2$)
                    \If{$r_1$}
                        \State $ isInvalid \gets true $
                        \State Continue
                    \Else
                        \State $g_3 \gets i_1$
                        \State $F \gets F \cup (g_3, g_2)$
                        \State $Q$.push($(g_3, g_2)$)
                        \State $Q$.push($lemma, counterExample$)
                        \State Continue
                    \EndIf
                \EndIf
                \State $g_1 \gets$ Generalize($m_1$, $T_k$, $\neg lemma$)
                \State $(r_2, i_2, n_2) \gets $ Reachable($n-k+1$, $n$, $g_1$)
                \If{$r_2$}
                    \State $(r_3, i_3, n_3) \gets $ Reachable($n+1$, $n_2 + k$, $g_1$)
                    \State $n_p \gets $ Min($n_p$, $n_3$)
                    \State $F \gets$ $F$ $\cup $ ($\neg counterExample$, $counterExample$)
                    \State $G \gets$ $G$ $\cup $ ($\neg counterExample$, $counterExample$)
                \Else
                    \State $g_3 \gets$ $i_2 \wedge lemma$
                    \State $F \gets$ $F$ $\cup $ ($g_3$, $counterExample$)
                    \State $F \gets$ $F$ $\setminus $ ($lemma$, $counterExample$)
                    \State $Q$.push(($g_3$, $counterExample$))
                \EndIf
                \State \Return ($F$, $G$, $n_p$, $isInvalid$)
            \EndWhile

        \end{algorithmic}
    \end{mdframed}
    \caption{Push procedure}\label{alg:3}
\end{figure}

\newpage
\section{PD-Kind procedure}

\noindent The main PDKind procedure shown in Algorithm \ref{alg:4}
checks if property \( P \) is invariant by
iteratively calling the \( Push \) procedure to find a k-inductive
strengthening of \( P \) for some $1 \leq k \leq n + 1$.
PDKind iteratively strengthens the invariant until it becomes inductive.
The strengthening
\( G \) is k-inductive, and if $F == G$, then no more strengthening is required, and \( P \) is proven to be invariant
,so we return \( SAFE \). If the \( Push \) procedure marks
\( isInvalid \) as \( True \), then a counterexample has been found, proving that the property is not invariant and we
return \( UNSAFE \). Otherwise, we update \( n \) and repeat the loop.


\begin{figure}[H]
\begin{mdframed}
\begin{algorithmic}[1]
\State \textbf{Input:} Initial states $I$, transition formula $T$, property $P$
\State \textbf{Output:} Return UNSAFE if P is invalid or SAFE when there is no inductive strengthening left

\State $n \gets 0$
    \State $F \gets (P, \neg P)$
    \While{true}
        \State $k \gets n + 1$
        \State $(F, G, n_p, isInvalid) \gets$ Push($F$, $n$, $k$)
        \If{$isInvalid$}
        \State \Return UNSAFE
        \EndIf
        \If{$F = G$}
        \State \Return SAFE
        \EndIf
        \State $n \gets n_p$
        \State $F \gets G$
    \EndWhile

\end{algorithmic}
\end{mdframed}
\caption{Main PD-Kind procedure}\label{alg:4}
\end{figure}


\chapter{Implementation}
% Zde bych popsal pouzity programovaci jazyk, doporucenou platformu.

% Nevim jestli tato cast nebude moc kratka, ale popsal bych jednotlive, jak reprezentuji ruzne datove struktury popsane v analyze.
\section{Data structures}

\subsection{Reachability frame}

\subsection{Induction frame}

% Zde bych popsal rozdeleni kodu na tridu, ktera se stara o Reachability check a na engine PDKIND
\section{Classes}

\subsection{Reachability class}

\subsection{PD-Kind engine}

\chapter{Experiments}

\chapter{Conclusion}
During this work, we introduced ourselves with the concept of the CHC framework. Then, we defined the concepts of transition systems and satisfiability modulo theories. After that, we described the structure of Golem and analyzed the best way to integrate a new engine into it. We then analyzed the entire PDKind algorithm and modified it to serve our needs.

With this knowledge, we implemented the PDKind engine itself. The goal was to create an engine that would 1) be well integrated into the Golem framework, 2) return correct answers, and 3) match the performance of other engines. To validate these goals, we experimentally compared our engine with other existing Golem engines. We managed to achieve results comparable to other engines, and in some cases, our engine was faster, although it did not reach the performance of the best-performing engine.

We believe that this engine presents a well-performing alternative to other engines, with its usefulness expected to grow with further improvements.


%%% Bibliography
%%% Bibliography (literature used as a source)
%%%
%%% We employ biblatex to construct the bibliography. It processes
%%% citations in the text (e.g., the \cite{...} macro) and looks up
%%% relevant entries in the bibliography.bib file.
%%%
%%% See also biblatex settings in thesis.tex.

%%% Generate the bibliography. Beware that if you cited no works,
%%% the empty list will be omitted completely.

% We let bibliography items stick out of the right margin a little
\def\bibfont{\hfuzz=2pt}

\printbibliography[heading=bibintoc]

%%% If case you prefer to write the bibliography manually (without biblatex),
%%% you can use the following. Please follow the ISO 690 standard and
%%% citation conventions of your field of research.

% \begin{thebibliography}{99}
%
% \bibitem{lamport94}
%   {\sc Lamport,} Leslie.
%   \emph{\LaTeX: A Document Preparation System}.
%   2nd edition.
%   Massachusetts: Addison Wesley, 1994.
%   ISBN 0-201-52983-1.
%
% \end{thebibliography}


%%% Figures used in the thesis (consider if this is needed)
\listoffigures

%%% Tables used in the thesis (consider if this is needed)
%%% In mathematical theses, it could be better to move the list of tables to the beginning of the thesis.
\listoftables

%%% Abbreviations used in the thesis, if any, including their explanation
%%% In mathematical theses, it could be better to move the list of abbreviations to the beginning of the thesis.
\chapwithtoc{List of Abbreviations}

%%% Doctoral theses must contain a list of author's publications
\ifx\ThesisType\TypePhD
\chapwithtoc{List of Publications}
\fi

%%% Attachments to the thesis, if any. Each attachment must be referred to
%%% at least once from the text of the thesis. Attachments are numbered.
%%%
%%% The printed version should preferably contain attachments, which can be
%%% read (additional tables and charts, supplementary text, examples of
%%% program output, etc.). The electronic version is more suited for attachments
%%% which will likely be used in an electronic form rather than read (program
%%% source code, data files, interactive charts, etc.). Electronic attachments
%%% should be uploaded to SIS. Allowed file formats are specified in provision
%%% of the rector no. 72/2017. Exceptions can be approved by faculty's coordinator.
\appendix
\chapter{Attachments}

\section{First Attachment}

\end{document}
